%  Математическая физика — конспект лекций
%  Компиляция: pdflatex main.tex  (дважды, из-за ссылок и оглавления)

\documentclass[12pt,a4paper]{article}

%  Преамбула конспекта «Математическая физика»
%  Компилировать pdfLaTeX (кодировка utf8, шрифты T2A).

% Latin Modern не имеет T2A-начертаний (нет t2almr.fd), поэтому кириллица идёт
% через cm-super: семейство cmr в кодировке T2A.
\usepackage[T2A]{fontenc}
\usepackage[utf8]{inputenc}
\usepackage[english,russian]{babel}

\usepackage{amsmath,amssymb,amsthm,mathtools}
\usepackage{geometry}
\geometry{a4paper,margin=2.5cm}

\usepackage{tikz}
\usetikzlibrary{arrows.meta,patterns,decorations.pathmorphing,decorations.pathreplacing,shapes.geometric,calc}
\usepackage{pgfplots}
\usepgfplotslibrary{groupplots,fillbetween}
\pgfplotsset{compat=1.18}

\usepackage[most]{tcolorbox}
\tcbuselibrary{breakable}

\usepackage{enumitem}
\usepackage{booktabs}
\usepackage{hyperref}
\hypersetup{
  colorlinks=true,
  linkcolor=blue!55!black,
  citecolor=blue!55!black,
  urlcolor=blue!55!black,
  pdftitle={Математическая физика --- конспект}
}

%  Окружения для теорем
\theoremstyle{plain}
\newtheorem{theorem}{Теорема}
\newtheorem{proposition}[theorem]{Утверждение}
\newtheorem{corollary}{Следствие}
\newtheorem*{corollary*}{Следствие}
\newtheorem{lemma}[theorem]{Лемма}
\newtheorem{criterion}{Признак}

\theoremstyle{definition}
\newtheorem{definition}{Определение}
\newtheorem{example}{Пример}
\newtheorem{exercise}{Упражнение}

\theoremstyle{remark}
\newtheorem*{remark}{Замечание}
\newtheorem*{reminder}{Напоминание}

\renewcommand{\qedsymbol}{$\blacksquare$}

%  Пометки конспекта
%    \begin{idea} ... \end{idea}  — пояснения, которых не было на лекции
%    \begin{warn} ... \end{warn} — неразборчиво / описка / требует сверки
\newtcolorbox{idea}[1][]{
  breakable,
  enhanced,
  colback=blue!4!white,
  colframe=blue!45!black,
  boxrule=0.4pt,
  left=6pt,right=6pt,top=5pt,bottom=5pt,
  title={Пояснение},
  fonttitle=\bfseries\small,
  coltitle=white,
  attach boxed title to top left={xshift=6pt,yshift=-2pt},
  boxed title style={colback=blue!45!black,boxrule=0pt,arc=1pt},
  #1
}

\newtcolorbox{warn}[1][]{
  breakable,
  enhanced,
  colback=orange!6!white,
  colframe=orange!65!black,
  boxrule=0.4pt,
  left=6pt,right=6pt,top=5pt,bottom=5pt,
  title={Сверить с записью},
  fonttitle=\bfseries\small,
  coltitle=white,
  attach boxed title to top left={xshift=6pt,yshift=-2pt},
  boxed title style={colback=orange!65!black,boxrule=0pt,arc=1pt},
  #1
}

% Отметка начала лекции
\newcommand{\lecture}[2]{%
  \bigskip
  \noindent\hfill\textit{\small #1\ifx\relax#2\relax\else, #2\fi}\\[-0.3em]
  \noindent\rule{\textwidth}{0.4pt}
  \medskip
}

%  Макросы
\newcommand{\R}{\mathbb{R}}
\newcommand{\N}{\mathbb{N}}
\newcommand{\Om}{\Omega}
\newcommand{\clOm}{\overline{\Omega}}
\newcommand{\dOm}{\partial\Omega}
\newcommand{\eps}{\varepsilon}
\newcommand{\dist}{\operatorname{dist}}
\newcommand{\supp}{\operatorname{supp}}
\newcommand{\esssup}{\operatorname*{ess\,sup}}
\newcommand{\Lp}[1]{L^{#1}}
\newcommand{\norm}[1]{\left\lVert #1 \right\rVert}
\newcommand{\abs}[1]{\left\lvert #1 \right\rvert}
\newcommand{\Lpn}[2]{\norm{#1}_{L^{#2}(\Omega)}}
\newcommand{\ball}[2]{B(#1,#2)}
\newcommand{\unif}{\rightrightarrows}
\newcommand{\Lloc}{L^1_{\mathrm{loc}}}
\newcommand{\Wp}[2]{W_{#1}^{#2}}
\newcommand{\Wpz}[2]{\mathring{W}_{#1}^{#2}}
\newcommand{\sgn}{\operatorname{sign}}
\newcommand{\dd}[2]{\frac{\partial #1}{\partial #2}}
\newcommand{\intR}{\int\limits_{-\infty}^{+\infty}}
\newcommand{\Bxat}[2]{B_{#1,#2}}

% Маркер «++ в кружке» из тетради: доказательство будет дописано позже
% по записям прошлых семестров.
\newcommand{\plusplus}{%
  \tikz[baseline=(pp.base)]{%
    \node[draw, ellipse, inner xsep=3pt, inner ysep=1pt, line width=0.4pt] (pp)
      {\footnotesize $+\,+$};}%
}

%  Вспомогательные функции для рисунков
%    sstep(t)      — гладкая «ступенька» 6t^5-15t^4+10t^3, склеенная с 0 и 1
%    urho(x,rho)   — схематическое усреднение ступеньки u = 1_{[0,1)}
%    bump(x,rho)   — ядро усреднения omega_rho в размерности n = 1
\pgfmathdeclarefunction{sstep}{1}{%
  \pgfmathparse{(#1<0) ? 0 : ((#1>1) ? 1 : (6*(#1)^5-15*(#1)^4+10*(#1)^3))}%
}
\pgfmathdeclarefunction{urho}{2}{%
  \pgfmathparse{sstep((#1+#2)/(2*#2)) - sstep((#1-1+#2)/(2*#2))}%
}
% У края носителя 1-(x/rho)^2 -> 0, и деление 1/(...) переполняет pgfmath
% (предел ~16384) ещё до вызова exp. Поэтому знаменатель зажат снизу через max:
% показатель никогда не превосходит 15 по модулю, exp(-15) ~ 3e-7 — на графике ноль.
\pgfmathdeclarefunction{bump}{2}{%
  \pgfmathparse{(abs(#1)>=0.999*#2) ? 0 : (2.2523/#2)*exp(-1/max(1-((#1)/#2)^2,0.0667))}%
}

%  Быстрая сборка: \showfiguresfalse выключает все TikZ-рисунки.
%  Нужно, когда бесплатный Overleaf не успевает в лимит времени.
\newif\ifshowfigures
\showfigurestrue

\newcommand{\inputfig}[1]{%
  \ifshowfigures
    \input{#1}%
  \else
    \begin{center}\fbox{\small рисунок выключен: \texttt{\detokenize{#1}}}\end{center}%
  \fi
}


\title{\textbf{Математическая физика}\\[0.3em]\large Конспект лекций}
\author{}
\date{}

\begin{document}

\maketitle
\thispagestyle{empty}

\begin{idea}[title={Обозначения в конспекте}]
Блок \textbf{«Пояснение»} — то, чего не было на лекции: интуиция, обоснование
пропущенных шагов, связи с другими темами.

\medskip
Блок \textbf{«Сверить с записью»} — места, где запись на доске или в тетради
неразборчива, содержит описку или требует уточнения.

\medskip
Маркер \plusplus{} вместо доказательства — так помечено в тетради; эти доказательства
дописываются позже по записям прошлых семестров.
\end{idea}

\tableofcontents
\newpage

% Нумерация: section = глава, subsection = параграф (§)
\renewcommand{\thesubsection}{\S\arabic{subsection}}

\section*{Напоминание: задача Коши для уравнения теплопроводности}
\addcontentsline{toc}{section}{Напоминание: уравнение теплопроводности}

Лекция началась с того, что было получено в прошлом семестре. Рассматривалась задача
Коши
\[
  \begin{cases}
    u_t - \alpha u_{xx} = f(x,t), & x \in \R,\ t > 0, \\[2pt]
    u|_{t=0} = \varphi(x),
  \end{cases}
\]
её решение записывается через фундаментальное решение $\Phi$:
\[
  u(x,t) = \intR \Phi(x-y,\,t)\,\varphi(y)\,dy
         + \int\limits_0^t \intR \Phi(x-y,\,t-\tau)\,f(y,\tau)\,dy\,d\tau ,
\]
\[
  \boxed{\ \Phi(x,t) = \frac{1}{2\sqrt{\pi \alpha t}}\;
         e^{-\frac{x^2}{4\alpha t}}\ }
\]

\inputfig{figures/fig-heat-kernel}

\begin{warn}
Рядом в тетради есть плохо читаемая приписка лектора — что-то вроде «как бы то
ни было, это всё приближение; ряд Фурье — это приближение, которое работает на
макроуровне». Смысл, видимо, в том, что само уравнение теплопроводности и его
решение описывают среду усреднённо, а не на уровне отдельных частиц. Стоит
уточнить формулировку.
\end{warn}

\begin{idea}
Структура ответа — свёртка начальных данных с $\Phi$ плюс интеграл Дюамеля по
источнику $f$ — повторится и для волнового уравнения: там роль $\Phi$ сыграет
полусумма сдвигов и интеграл по характеристическому треугольнику.

\medskip
Два свойства $\Phi$ объясняют картинку: $\intR \Phi(x,t)\,dx = 1$ при любом $t>0$
(«тепло никуда не девается»), а высота пика равна
$\Phi(0,t) = \frac{1}{2\sqrt{\pi\alpha t}} \to \infty$ при $t \to 0+$. То есть
$\Phi(\cdot,t) \to \delta$ при $t \to 0+$. При этом $\Phi(x,t) > 0$ сразу для всех
$x$: возмущение в уравнении теплопроводности распространяется мгновенно. Дальше
будет видно, что у волнового уравнения всё наоборот — конечная скорость $a$.
\end{idea}


\newpage
\section{Волновое уравнение}
\subsection{Колебание бесконечной струны}

\subsubsection*{Свободные колебания: задача (1)}

\stepcounter{equation}% уникальный якорь hyperref: \tag сам счётчик не двигает
\begin{equation}\tag{1}\label{eq:inf-free}
  \begin{cases}
    u_{tt} - a^2 u_{xx} = 0, & x \in \R,\ t > 0, \\[2pt]
    u|_{t=0} = \varphi(x), \quad u_t|_{t=0} = \psi(x).
  \end{cases}
\end{equation}

\begin{warn}
На доске в первой записи коэффициент выглядит как $v^2$, во всех последующих —
как $a^2$; принято $a^2$. Там же в условии написано «$t = \infty$» — очевидная
описка вместо $t > 0$.
\end{warn}

Решение задачи~\eqref{eq:inf-free} получается заменой через характеристики —
той же самой $\xi = x - at$, $\eta = x + at$, что подробно разобрана ниже для
задачи~\eqref{eq:inf-forced}:
\[
  u(x,t) = \frac12\Big[\varphi(x-at) + \varphi(x+at)\Big]
         + \frac{1}{2a}\int\limits_{x-at}^{x+at} \psi(y)\, dy
\]
\begin{center}
\textbf{— формула Даламбера—Эйлера.}
\end{center}

\begin{warn}
На доске пределы интеграла от $\psi$ в этом месте записаны как $-\infty$ и
$+\infty$. Из итоговой формулы Даламбера (см.~\eqref{eq:dalembert}) видно, что
пределы — $x-at$ и $x+at$.
\end{warn}

\begin{idea}
Откуда берётся эта формула: оператор факторизуется,
\[
  \partial_{tt} - a^2 \partial_{xx}
  = (\partial_t - a\partial_x)(\partial_t + a\partial_x),
\]
поэтому общее решение однородного уравнения — сумма двух бегущих волн
$u = F(x-at) + G(x+at)$. Подстановка начальных условий даёт
\[
  F + G = \varphi, \qquad a\,(G' - F') = \psi,
\]
откуда и получаются полусумма сдвигов и интеграл от $\psi$.
\end{idea}

\begin{example}[начальное отклонение, нулевая скорость]
$\psi \equiv 0$, $\varphi$ — «треугольник». Тогда
$u(x,t) = \tfrac12[\varphi(x-at) + \varphi(x+at)]$: начальный горб распадается на
две волны половинной высоты, уходящие налево и направо со скоростью $a$.
\end{example}

\inputfig{figures/fig-wave-split}

\begin{example}[нулевое отклонение, точечный удар]
$\varphi \equiv 0$, $\psi(x) = \delta(x)$. Тогда
\[
  u(x,t) = \frac{1}{2a}\int\limits_{x-at}^{x+at} \delta(y)\, dy
  = \begin{cases}
      \dfrac{1}{2a}, & |x| < at, \\[6pt]
      0, & |x| > at,
    \end{cases}
\]
то есть «ступенька» высоты $\tfrac{1}{2a}$, края которой — точки $-at$ и $at$ —
разбегаются со скоростью $a$. Разрывный профиль — это, разумеется, просто
математическая модель: реальная струна таких углов не имеет.
\end{example}

\inputfig{figures/fig-delta-box}

\begin{idea}
Разница с уравнением теплопроводности видна сразу: там точечный источник мгновенно
«размазывается» по всей прямой (рис.~\ref{fig:heat-kernel}), здесь возмущение строго
заперто в интервале $|x| < at$. Это и есть конечность скорости распространения.
\end{idea}

\subsubsection*{Вынужденные колебания при нулевых данных: задача (2)}

\stepcounter{equation}
\begin{equation}\tag{2}\label{eq:inf-forced}
  \begin{cases}
    u_{tt} - a^2 u_{xx} = f(x,t), & x \in \R,\ t > 0, \\[2pt]
    u|_{t=0} = 0, \quad u_t|_{t=0} = 0 .
  \end{cases}
\end{equation}

Заметим сразу: $u|_{t=0} = 0 \Longrightarrow u_x|_{t=0} = 0$ (дифференцируем
тождественный нуль вдоль прямой $t=0$).

\paragraph{Характеристические переменные.} Положим
\[
  \begin{cases} \xi = x - at, \\ \eta = x + at, \end{cases}
  \qquad\Longleftrightarrow\qquad
  \begin{cases} x = \tfrac12(\xi + \eta), \\[2pt] t = \tfrac{1}{2a}(\eta - \xi). \end{cases}
\]
В этих переменных
\[
  u_{tt} - a^2 u_{xx} = -4a^2 u_{\xi\eta},
  \qquad\text{так что}\qquad
  -4a^2 u_{\xi\eta} = \widetilde f(\xi,\eta),
\]
где $\widetilde f(\xi,\eta) = f\!\left(\tfrac12(\xi+\eta),\ \tfrac{1}{2a}(\eta-\xi)\right)$.

Прямая $t = 0$ переходит в прямую $\xi = \eta$. Начальные условия дают на ней
\[
  u|_{\xi=\eta} = 0, \qquad u_\xi|_{\xi=\eta} = 0, \qquad u_\eta|_{\xi=\eta} = 0 .
\]

\begin{idea}
Последние два равенства — это пересчёт условий $u_x|_{t=0} = 0$ и $u_t|_{t=0}=0$:
\[
  u_x = u_\xi + u_\eta, \qquad u_t = a\,(u_\eta - u_\xi),
\]
и если обе левые части равны нулю на прямой $t=0$, то там $u_\xi = u_\eta = 0$.
\end{idea}

\paragraph{Интегрирование по характеристическому треугольнику.}
Зафиксируем точку $(\xi_0, \eta_0)$ и возьмём треугольник
\[
  D = \{\, (\xi,\eta) : \ \xi_0 \le \xi \le \eta_0, \ \ \xi \le \eta \le \eta_0 \,\}.
\]

\inputfig{figures/fig-char-xieta}

Тогда
\begin{align*}
  \iint\limits_D u_{\xi\eta} \, d\xi\, d\eta
  &= \int\limits_{\xi_0}^{\eta_0} \!\! \int\limits_{\xi}^{\eta_0}
     u_{\xi\eta}\, d\eta\, d\xi
   = \int\limits_{\xi_0}^{\eta_0} u_\xi \Big|_{\eta=\xi}^{\eta=\eta_0} d\xi
   = \int\limits_{\xi_0}^{\eta_0} u_\xi(\xi, \eta_0)\, d\xi \\[4pt]
  &= u(\xi,\eta_0)\Big|_{\xi=\xi_0}^{\xi=\eta_0}
   = -\,u(\xi_0, \eta_0),
\end{align*}
поскольку $u_\xi|_{\eta=\xi} = 0$, а $u(\eta_0,\eta_0) = 0$ (эта точка лежит на
прямой $\xi=\eta$). С другой стороны $u_{\xi\eta} = -\tfrac{1}{4a^2}\widetilde f$,
значит
\[
  \boxed{\ u(\xi_0,\eta_0) = \frac{1}{4a^2} \iint\limits_D \widetilde f(\xi,\eta)\,
  d\xi\, d\eta \ }
\]

\paragraph{Возврат к переменным $(x,t)$.} Якобиан замены:
\[
  J = \det \begin{pmatrix} 1 & -a \\ 1 & a \end{pmatrix} = 2a ,
  \qquad d\xi\, d\eta = 2a \, dy \, d\tau .
\]
Поэтому
\[
  u(x_0,t_0) = \frac{1}{4a^2} \iint\limits_{D} 2a\, f(y,\tau)\, dy\, d\tau
             = \frac{1}{2a} \iint\limits_{D} f(y,\tau)\, dy\, d\tau ,
\]
где теперь $D$ — характеристический треугольник в плоскости $(x,t)$ с вершиной
$(x_0,t_0)$ и боковыми сторонами $|x - x_0| = a(t_0 - t)$:
\[
  u(x_0,t_0) = \frac{1}{2a} \int\limits_0^{t_0}
  \int\limits_{x_0 - a(t_0-\tau)}^{\,x_0 + a(t_0-\tau)} f(y,\tau)\, dy\, d\tau .
\]

\inputfig{figures/fig-char-triangle}

\begin{warn}
И на доске, и в тетради пределы внутреннего интеграла записаны короче — как
$x_0 - at$ и $x_0 + at$ (а ниже, в формуле Даламбера, как $x - a\tau$ и
$x + a\tau$). Поскольку по внешней переменной идёт интегрирование, формально
это неверно: по рисунку характеристического треугольника срез на уровне $\tau$
имеет концы $x_0 \mp a(t_0 - \tau)$, и именно так записано выше. Судя по тому,
что сокращение повторяется во всех записях, это устойчивая скоропись лектора,
а не другая формула.
\end{warn}

\subsubsection*{Общая задача (3) и формула Даламбера}

\stepcounter{equation}
\begin{equation}\tag{3}\label{eq:inf-full}
  \begin{cases}
    u_{tt} - a^2 u_{xx} = f(x,t), & x \in \R,\ t > 0, \\[2pt]
    u|_{t=0} = \varphi(x), \quad u_t|_{t=0} = \psi(x) .
  \end{cases}
\end{equation}

По линейности задача~\eqref{eq:inf-full} есть сумма задач~\eqref{eq:inf-free}
и~\eqref{eq:inf-forced}, поэтому

\stepcounter{equation}
\begin{equation}\tag{Д}\label{eq:dalembert}
  \boxed{
  \begin{aligned}
    u(x,t) &= \frac12\Big[\varphi(x-at) + \varphi(x+at)\Big]
            + \frac{1}{2a}\int\limits_{x-at}^{x+at} \psi(y)\, dy \\[4pt]
           &\quad + \frac{1}{2a} \int\limits_0^t
             \int\limits_{x-a(t-\tau)}^{\,x+a(t-\tau)} f(y,\tau)\, dy\, d\tau
  \end{aligned}}
\end{equation}

\begin{center}
\textbf{— формула Даламбера.}
\end{center}

\subsubsection*{Компактная запись}

Обозначим интервал
\[
  \Bxat{x}{at} = (x - at,\ x + at).
\]
Заметим, что
\[
  \frac12\Big[\varphi(x-at) + \varphi(x+at)\Big]
  = \frac{1}{2a}\,\dd{}{t} \int\limits_{x-at}^{x+at} \varphi(y)\, dy
  = \frac{1}{2a}\,\dd{}{t} \int\limits_{\Bxat{x}{at}} \varphi(y)\, dy ,
\]
поскольку дифференцирование интеграла по пределам даёт
$a\varphi(x+at) + a\varphi(x-at)$. Поэтому формула Даламбера записывается
единообразно:
\[
  \boxed{\ u(x,t) = \frac{1}{2a}\left[
    \dd{}{t} \int\limits_{\Bxat{x}{at}} \varphi(y)\, dy
    + \int\limits_{\Bxat{x}{at}} \psi(y)\, dy
    + \int\limits_0^t \int\limits_{\Bxat{x}{a(t-\tau)}} f(y,\tau)\, dy\, d\tau
  \right]\ }
\]

\begin{warn}
В записи интервала на доске стояло $\Bxat{x}{at} = (x+at,\, x+at)$; в тетради
записано правильно, $(x-at,\, x+at)$, — значит на доске просто описка.
А вот в третьем слагаемом и доска, и тетрадь пишут $\Bxat{x}{a\tau}$: это то же
сокращение, что и выше, по характеристическому треугольнику там
$\Bxat{x}{a(t-\tau)}$.
\end{warn}

\begin{idea}
Запись через $\Bxat{x}{at}$ не косметическая: в размерности $n = 3$ (см.~\S3)
формула будет выглядеть ровно так же, только $\Bxat{x}{at}$ станет шаром радиуса
$at$, а интегралы — средними по его границе. Одномерный интервал — это «шар» при
$n=1$.
\end{idea}

\subsubsection*{Существование классического решения}

\begin{theorem}\label{thm:dalembert-exist}
Пусть
\[
  \varphi \in C^2(\R), \qquad \psi \in C^1(\R), \qquad f,\ f_x \in C(\R) .
\]
Тогда классическое решение задачи~\eqref{eq:inf-full} существует и выражается
формулой Даламбера~\eqref{eq:dalembert}.
\end{theorem}

\begin{proof}
Уже доказано: проведённый выше вывод и есть доказательство. Остаётся заметить,
что при указанной гладкости формула Даламбера дважды непрерывно
дифференцируема, то есть даёт \emph{классическое} решение.
\end{proof}

\begin{idea}
Требования по гладкости ровно такие, чтобы формула Даламбера после двукратного
дифференцирования дала непрерывные $u_{tt}$ и $u_{xx}$.

\medskip
$\varphi \in C^2$ — входит без интегрирования, поэтому теряет обе производные
целиком.

\medskip
$\psi \in C^1$ — стоит под интегралом с пределами, зависящими от $x$ и $t$.
Дифференцирование по пределам «съедает» одну производную: первая получается
бесплатно, на вторую нужна гладкость самой $\psi$.

\medskip
$f, f_x \in C$ — $f$ стоит под двойным интегралом, поэтому по $t$ запаса
хватает; но при дифференцировании по $x$ один раз интеграл снимается
пределами, а для второго раза нужна уже непрерывность $f_x$. Отсюда и
несимметричность условия по $x$ и $t$.
\end{idea}

\subsubsection*{Область зависимости}

\inputfig{figures/fig-dependence}

\begin{warn}
В тетради треугольник подписан «область влияния \textbf{на точку} $(x,t)$» — с
этим уточнением всё согласуется: речь именно о том, что влияет \emph{на} точку,
то есть об области зависимости. На доске слова «на точку» потерялись, оттого и
возникала путаница с конусом \emph{над} точкой. Для справки, стандартная
терминология:
треугольник \emph{под} точкой $(x,t)$ — то, от чего значение $u(x,t)$ зависит, —
называется \textbf{областью зависимости}, а областью влияния точки называют конус
\emph{над} ней.
\end{warn}

\begin{idea}
Содержательно: значение $u(x,t)$ определяется данными $\varphi, \psi$ только на
интервале $\Bxat{x}{at}$ и значениями $f$ только внутри характеристического
треугольника. Изменим данные вне этого треугольника — значение $u(x,t)$ не
изменится.

\medskip
Это и есть \textbf{конечная скорость распространения}: события в точке $(x,t)$
могут не зависеть от более ранних событий, если те произошли достаточно далеко.
У уравнения теплопроводности скорость распространения бесконечна — там
возмущение из любой точки мгновенно сказывается всюду
(ср.~рис.~\ref{fig:heat-kernel}).
\end{idea}

\subsubsection*{Непрерывная зависимость от данных}

Пусть двум наборам данных отвечают два решения:
\[
  \varphi_1, \psi_1, f_1 \ \rightsquigarrow\ u_1,
  \qquad
  \varphi_2, \psi_2, f_2 \ \rightsquigarrow\ u_2 .
\]
Введём
\[
  \delta_\varphi := \sup_x \abs{\varphi_1(x) - \varphi_2(x)}, \qquad
  \delta_\psi := \sup_x \abs{\psi_1(x) - \psi_2(x)},
\]
\[
  \delta_{f,T} := \sup_{\substack{x \in \R \\ t \in [0,T]}}
    \abs{f_1(x,t) - f_2(x,t)} .
\]

\begin{theorem}
Классическое решение задачи~\eqref{eq:inf-full} непрерывно зависит от
$f, \varphi, \psi$ в том смысле, что для любого $T > 0$ выполняется неравенство
\[
  \abs{u_1(x,t) - u_2(x,t)} \;\le\; \delta_\varphi + \delta_\psi\, T
    + \delta_{f,T}\,\frac{T^2}{2}
  \qquad \forall\, x\in\R,\ t \in [0,T] .
\]
\end{theorem}

\begin{proof}
Оставлено как домашнее задание; сама выкладка разобрана в пояснении ниже.
\end{proof}

Единственность здесь не доказывается: она будет получена позже и сразу для
$n$-мерного случая (\S4).

\begin{idea}
Оценка читается прямо из формулы Даламбера: разность решений — это та же формула,
применённая к разностям данных. Полусумма сдвигов даёт $\delta_\varphi$; интеграл
от $\psi$ по интервалу длины $2at$ даёт
$\frac{1}{2a}\cdot 2at \cdot \delta_\psi = \delta_\psi\, t$; двойной интеграл по
треугольнику площади $a t^2$ даёт
$\frac{1}{2a}\cdot a t^2 \cdot \delta_{f,T} = \delta_{f,T}\,\frac{t^2}{2}$. Затем
$t$ оценивается через $T$.

\medskip
На доске эта выкладка сначала записана с $t$, а в формулировке теоремы — уже с
$T$; это одно и то же неравенство.

\medskip
Вместе с теоремой~\ref{thm:dalembert-exist} и единственностью (она следует из
самой формулы: решение задаётся данными однозначно) это даёт
\textbf{корректность по Адамару} задачи~\eqref{eq:inf-full}: существование,
единственность, непрерывная зависимость от данных.
\end{idea}


\newpage
\subsection{Свободные колебания полубесконечной струны}

\stepcounter{equation}
\begin{equation}\tag{1}\label{eq:semi-dirichlet}
  \begin{cases}
    u_{tt} - a^2 u_{xx} = 0, & x > 0,\ t > 0, \\[2pt]
    u|_{t=0} = \varphi(x), \quad u_t|_{t=0} = \psi(x), \\[2pt]
    u|_{x=0} = 0 .
  \end{cases}
\end{equation}

Теперь струна занимает полупрямую, а её левый конец закреплён. Нумерация задач в
этом параграфе начинается заново.

\subsubsection*{Метод нечётного продолжения}

Продолжим данные на всю ось нечётно:
\[
  \Phi(x) =
  \begin{cases}
    \varphi(x), & x \ge 0, \\
    -\varphi(-x), & x \le 0,
  \end{cases}
  \qquad\qquad
  \Psi(x) =
  \begin{cases}
    \psi(x), & x \ge 0, \\
    -\psi(-x), & x \le 0 .
  \end{cases}
\]

\inputfig{figures/fig-odd-extension}

Решим на всей прямой задачу с данными $\Phi, \Psi$ по формуле Даламбера:
\[
  u(x,t) = \frac12\Big[\Phi(x-at) + \Phi(x+at)\Big]
         + \frac{1}{2a}\int\limits_{x-at}^{x+at} \Psi(y)\, dy .
\]

\begin{idea}
Ключ: если $\Phi$ и $\Psi$ нечётны, то $u|_{x=0} = 0$ автоматически. Действительно,
при $x=0$
\[
  \frac12\big[\Phi(-at) + \Phi(at)\big] = 0
  \quad\text{(нечётность $\Phi$)},
\]
\[
  \int\limits_{-at}^{at} \Psi(y)\, dy = 0
  \quad\text{(интеграл нечётной функции по симметричному отрезку)} .
\]
Значит построенная функция удовлетворяет и уравнению, и начальным данным (при
$x>0$ они совпадают с исходными), и краевому условию. Условие Дирихле
$u|_{x=0}=0$ «встроено» в симметрию данных.
\end{idea}

\subsubsection*{Явный вид решения}

Если $x - at \ge 0$, то все аргументы попадают на полупрямую $x\ge 0$ и формула
не отличается от случая бесконечной струны. Если же
\[
  x - at < 0,
\]
то работает нечётность:
\[
  \Phi(x-at) = -\varphi(at-x),
  \qquad
  \int\limits_{x-at}^{at-x} \Psi(y)\, dy = 0 ,
\]
поэтому интеграл от $\Psi$ сворачивается до интеграла по отрезку
$[\,at-x,\ x+at\,]$ от исходной $\psi$. Итог:
\[
  \boxed{\
  u(x,t) =
  \begin{cases}
    \dfrac12\Big[\varphi(x-at) + \varphi(x+at)\Big]
      + \dfrac{1}{2a} \displaystyle\int\limits_{x-at}^{x+at} \psi(y)\, dy,
      & x \ge at, \\[18pt]
    \dfrac12\Big[\varphi(x+at) - \varphi(at-x)\Big]
      + \dfrac{1}{2a} \displaystyle\int\limits_{at-x}^{x+at} \psi(y)\, dy,
      & x < at .
  \end{cases}\ }
\]

\begin{idea}
Физический смысл второй строки: при $x < at$ до точки $x$ уже успела дойти волна,
отражённая от закреплённого конца. Слагаемое $-\varphi(at-x)$ — это как раз
отражённая волна, и знак «минус» означает отражение \emph{с переворотом}: у
закреплённого конца горб возвращается впадиной.

\medskip
Прямая $x = at$ в плоскости $(x,t)$ разделяет две зоны: ниже неё отражение ещё не
пришло, выше — уже пришло.
\end{idea}

\subsubsection*{Упражнения}

\begin{exercise}
Решите задачу (свободный конец вместо закреплённого; $\varphi, \psi$ заданы на
полуоси):
\stepcounter{equation}
\begin{equation}\tag{2}\label{eq:semi-neumann}
  \begin{cases}
    u_{tt} - a^2 u_{xx} = 0, & x > 0,\ t > 0, \\[2pt]
    u|_{t=0} = \varphi(x), \quad u_t|_{t=0} = \psi(x), \\[2pt]
    u_x|_{x=0} = 0 .
  \end{cases}
\end{equation}
\end{exercise}

\begin{idea}
Подсказка: условие $u_x|_{x=0} = 0$ (Нейман) получается тем же приёмом, но с
\textbf{чётным} продолжением данных. Отражение тогда идёт без переворота знака.
\end{idea}

\begin{exercise}
Изобразите решения задач~\eqref{eq:semi-dirichlet} и~\eqref{eq:semi-neumann} со
следующими начальными условиями:
\begin{enumerate}[label=(\alph*), leftmargin=2em]
  \item $\varphi(x)$ — «треугольник» с вершиной в точке $x_0 > 0$, $\psi \equiv 0$;
  \item $\varphi \equiv 0$, $\psi(x) = \delta(x - x_0)$.
\end{enumerate}
\end{exercise}

\begin{exercise}
Решите задачи для уравнения теплопроводности на полупрямой:
\[
  \begin{cases}
    u_t - \alpha u_{xx} = 0, & x > 0,\ t > 0, \\
    u|_{t=0} = \varphi(x), \\
    u|_{x=0} = 0 ;
  \end{cases}
  \qquad\qquad
  \begin{cases}
    u_t - \alpha u_{xx} = 0, & x > 0,\ t > 0, \\
    u|_{t=0} = \varphi(x), \\
    u_x|_{x=0} = 0 .
  \end{cases}
\]
\end{exercise}

\begin{idea}
Тот же метод отражения работает и здесь: нечётное продолжение для условия
Дирихле, чётное — для Неймана, а затем свёртка с $\Phi$ из напоминания в начале
конспекта.
\end{idea}


\newpage
\subsection{Колебание в трёхмерном пространстве}

\stepcounter{equation}
\begin{equation}\tag{1}\label{eq:wave3-free}
  \begin{cases}
    u_{tt} - a^2 \Delta u = 0, & x \in \R^3,\ t > 0, \\[2pt]
    u|_{t=0} = \varphi(x), \quad u_t|_{t=0} = \psi(x) .
  \end{cases}
\end{equation}

\subsubsection*{Сферические средние}

\begin{definition}
\textbf{Сферическим средним} функции $u$ называется
\[
  \bar u(x,t,z) := \frac{1}{4\pi z^2} \int\limits_{S_{x,z}} u(y,t)\, ds
  = \frac{1}{4\pi z^2} \int\limits_0^{2\pi} \int\limits_0^{\pi}
    u\, z^2 \sin\theta \, d\theta\, d\varphi ,
\]
где $S_{x,z}$ — сфера радиуса $z$ с центром в $x$.
\end{definition}

\begin{warn}
На доске радиальная переменная обозначена буквой, похожей на $z$; это именно
радиус, а не декартова координата. Вторая радиальная переменная, по которой идёт
интегрирование внутри шара, обозначена $\rho$. Здесь сохранены обозначения доски:
$z$ — радиус внешней сферы, $\rho \in [0,z]$ — переменная интегрирования.
\end{warn}

\inputfig{figures/fig-spherical-mean}

\begin{idea}
Множитель $4\pi z^2$ — это площадь сферы $S_{x,z}$, так что $\bar u$ —
действительно среднее по сфере, а не просто интеграл. Отсюда сразу
\[
  \bar u(x,t,z) \xrightarrow[z \to 0]{} u(x,t)
\]
для непрерывной $u$: среднее по стягивающейся сфере стремится к значению в
центре. Именно этот предельный переход в конце даст формулу Кирхгофа.
\end{idea}

\subsubsection*{Уравнение для сферического среднего}

Проинтегрируем уравнение по шару $B_{x,z}$. Левая часть:
\[
  \int\limits_{B_{x,z}} u_{tt}\, dx
  = \int\limits_0^{z} \int\limits_{S_{x,\rho}} u_{tt}\, ds\, d\rho
  = \int\limits_0^{z} 4\pi \rho^2\, \bar u_{tt}\, d\rho .
\]
Правая часть — по первой формуле Грина (она же формула
Гаусса—Остроградского для $\operatorname{div}\nabla u$), с учётом того, что
внешняя нормаль к сфере направлена по радиусу, то есть
$\partial/\partial n = \partial/\partial z$; затем переходим от сферы к
угловым переменным:
\[
  \int\limits_{B_{x,z}} \Delta u\, dx
  = \int\limits_{S_{x,z}} \dd{u}{n}\, ds
  = \int\limits_{S_{x,z}} \dd{u}{z}\, ds
  = \int\limits_0^{2\pi} \int\limits_0^{\pi} \dd{u}{z}\, z^2 \sin\theta\,
    d\theta\, d\varphi
  = z^2 \dd{}{z}\big(4\pi \bar u\big) .
\]
Приравнивая (уравнение даёт $u_{tt} = a^2 \Delta u$) и сокращая на $4\pi$:
\[
  \int\limits_0^{z} \rho^2\, \bar u_{tt}\, d\rho
    - a^2 z^2 \dd{\bar u}{z} = 0 .
\]
Дифференцируем по $z$:
\[
  z^2 \bar u_{tt} - a^2 \dd{}{z}\!\left(z^2 \dd{\bar u}{z}\right) = 0
  \qquad\Longleftrightarrow\qquad
  \bar u_{tt} - a^2 \frac{1}{z^2} \dd{}{z}\!\left(z^2 \dd{\bar u}{z}\right) = 0 .
\]

\begin{idea}
Это уравнение Эйлера—Пуассона—Дарбу: радиальная часть волнового оператора в
$\R^3$. Главное наблюдение следующего шага — в размерности $3$ (и только в
нечётных) оно сводится к \emph{одномерному} волновому уравнению заменой
$v = z\bar u$.
\end{idea}

\subsubsection*{Сведение к полубесконечной струне}

Положим
\[
  \bar u = \frac{v}{z}, \qquad\text{то есть}\qquad v(z,t) = z\,\bar u(x,t,z) .
\]
Тогда
\[
  \begin{cases}
    v_{tt} - a^2 v_{zz} = 0, & z \in (0,\infty), \\[2pt]
    v|_{t=0} = z\,\bar\varphi =: v_0(z), \quad
    v_t|_{t=0} = z\,\bar\psi =: v_1(z), \\[2pt]
    v|_{z=0} = 0 .
  \end{cases}
\]

\begin{idea}
Это ровно задача~\eqref{eq:semi-dirichlet} из \S2: свободные колебания
полубесконечной струны с закреплённым концом. Условие $v|_{z=0} = 0$ не
навязано извне — оно следует из определения $v = z\bar u$ при ограниченной
$\bar u$.

\medskip
Проверка замены: прямое вычисление даёт
$\frac{1}{z^2}\big(z^2 \bar u_z\big)_z = \frac{1}{z}\,(z\bar u)_{zz}$, откуда
после умножения уравнения на $z$ и получается $v_{tt} = a^2 v_{zz}$.
\end{idea}

Здесь $\bar\varphi$ и $\bar\psi$ — сферические средние начальных данных:
\[
  \bar\varphi(x,z) = \frac{1}{4\pi z^2} \int\limits_{S_{x,z}} \varphi(y)\, ds,
  \qquad
  \bar\psi(x,z) = \frac{1}{4\pi z^2} \int\limits_{S_{x,z}} \psi(y)\, ds .
\]

\subsubsection*{Предельный переход и формула Кирхгофа}

Искомое значение восстанавливается как предел среднего:
\[
  u(x,t) = \lim_{z \to 0} \bar u = \lim_{z \to 0} \frac{v(z,t)}{z}
         = \dd{v}{z}\bigg|_{z=0} ,
\]
по правилу Лопиталя: $v|_{z=0} = 0$, так что это неопределённость
$\tfrac{0}{0}$, и предел равен производной $v$ по $z$ в нуле.

Так как $z \to 0$, при фиксированном $t > 0$ мы попадаем в случай $z - at < 0$,
то есть во вторую строку ответа из \S2:
\[
  v = \frac12\Big[v_0(z+at) - v_0(at-z)\Big]
    + \frac{1}{2a}\int\limits_{at-z}^{z+at} v_1(\rho)\, d\rho .
\]
Дифференцируем по $z$ и полагаем $z = 0$:
\[
  \dd{v}{z}\bigg|_{z=0}
  = \frac12\Big[v_0'(at) + v_0'(at)\Big]
    + \frac{1}{2a}\Big[v_1(at) + v_1(at)\Big]
  = v_0'(at) + \frac1a\, v_1(at) .
\]

Остаётся подставить $v_0$ и $v_1$:
\[
  v_0(z) = z\,\bar\varphi = \frac{z}{4\pi z^2} \int\limits_{S_{x,z}} \varphi\, ds
         = \frac{1}{4\pi z} \int\limits_{S_{x,z}} \varphi\, ds ,
\]
\[
  \frac1a\, v_1\big|_{z=at} = \frac1a \big(z\,\bar\psi\big)\Big|_{z=at}
  = \frac1a \cdot \frac{at}{4\pi a^2 t^2} \int\limits_{S_{x,at}} \psi\, ds .
\]
Производная $v_0'$ в точке $z = at$ переписывается через $\partial/\partial t$:
\[
  v_0'(at) = \frac1a\, \dd{}{t}\Big[v_0(at)\Big] ,
\]
так что
\[
  u(x,t) = v_0'(at) + \frac1a v_1(at)
         = \frac1a\left[\dd{}{t} v_0(at) + v_1(at)\right] .
\]

Итого:

\stepcounter{equation}
\begin{equation}\tag{К}\label{eq:kirchhoff-free}
  \boxed{\ u(x,t) = \frac{1}{4\pi a^2}\left[
    \dd{}{t}\left(\frac1t \int\limits_{S_{x,at}} \varphi(y)\, ds_y\right)
    + \frac1t \int\limits_{S_{x,at}} \psi(y)\, ds_y
  \right]\ }
\end{equation}

\begin{center}
\textbf{— формула Кирхгофа} (однородное уравнение).
\end{center}

\begin{idea}
Сравните с компактной записью Даламбера из \S1: структура та же — производная по
$t$ от интеграла по «шару» от $\varphi$ плюс интеграл от $\psi$. Разница
принципиальная: при $n=1$ интегрирование шло по \emph{отрезку}
$\Bxat{x}{at}$, а при $n=3$ — по \emph{сфере} $S_{x,at}$, то есть только по
границе. Отсюда и принцип Гюйгенса (см. ниже).
\end{idea}

\subsubsection*{Неоднородное уравнение: принцип Дюамеля}

Теперь наоборот: однородные начальные условия и неоднородное уравнение. Здесь
\[
  \Box u = u_{tt} - a^2 \Delta u
\]
— \textbf{оператор Даламбера}.

\stepcounter{equation}
\begin{equation}\tag{2}\label{eq:wave3-forced}
  \begin{cases}
    \Box u \equiv u_{tt} - a^2 \Delta u = f(x,t), & x \in \R^3,\ t > 0, \\[2pt]
    u|_{t=0} = 0, \quad u_t|_{t=0} = 0 .
  \end{cases}
\end{equation}

Представим правую часть как суперпозицию мгновенных импульсов:
\[
  f(x,t) = \int\limits_0^{\infty} f(x,\tau)\, \delta(\tau - t)\, d\tau ,
\]
и рассмотрим вспомогательную задачу на один импульс:
\[
  \begin{cases}
    v_{tt} - a^2 \Delta v = f(x,\tau)\,\delta(t - \tau), & t > 0, \\[2pt]
    v|_{t=0} = 0, \quad v_t|_{t=0} = 0 .
  \end{cases}
\]

\paragraph{Что делает импульс.} Проинтегрируем уравнение по
$t \in [\tau-\eps,\, \tau+\eps]$ и устремим $\eps \to 0$:
\[
  \int\limits_{\tau-\eps}^{\tau+\eps} v_{tt}\, dt
  = v_t \Big|_{\tau-\eps}^{\tau+\eps}
  \xrightarrow[\eps \to 0]{} v_t\big|_{t=\tau} ,
  \qquad
  \int\limits_{\tau-\eps}^{\tau+\eps} \Delta v\, dt
  \xrightarrow[\eps \to 0]{} 0 ,
\]
\begin{idea}
Почему в первом пределе остаётся только верхний предел: при $t < \tau$ правая
часть тождественно равна нулю, а начальные данные нулевые, поэтому по теореме
единственности (\S4) там $v \equiv 0$. Значит и $v_t(\tau - \eps) \to 0$.
Во втором пределе всё проще: $\Delta v$ — хорошая функция, а промежуток
интегрирования стягивается в точку.
\end{idea}
\[
  \int\limits_{\tau-\eps}^{\tau+\eps} f(x,\tau)\,\delta(\tau - t)\, dt
  = f(x,\tau) .
\]
Значит $\delta$-импульс не меняет саму $v$, но мгновенно задаёт ей скорость:
\[
  \begin{cases}
    v_{tt} - a^2 \Delta v = 0, & t > \tau, \\[2pt]
    v|_{t=\tau} = 0, \quad v_t|_{t=\tau} = f(x,\tau) .
  \end{cases}
\]

\paragraph{Суперпозиция.} Положим
\[
  u(x,t) = \int\limits_0^{\infty} v(x,t,\tau)\, d\tau
         = \int\limits_0^{t} v\, d\tau
\]
(при $\tau > t$ решение $v$ ещё тождественно равно нулю).

\begin{proposition}
Так построенная $u$ решает задачу~\eqref{eq:wave3-forced}.
\end{proposition}

\begin{proof}
Начальные условия: $u|_{t=0} = 0$ очевидно, и
\[
  u_t = \underbrace{v\big|_{t=\tau}}_{=\,0} + \int\limits_0^t v_t\, d\tau
      = \int\limits_0^t v_t\, d\tau
  \qquad\Longrightarrow\qquad u_t|_{t=0} = 0 .
\]
Далее
\[
  u_{tt} = v_t\big|_{t=\tau} + \int\limits_0^t v_{tt}\, d\tau
         = f(x,t) + \int\limits_0^t v_{tt}\, d\tau ,
\]
и поэтому
\[
  u_{tt} - a^2 \Delta u
  = f(x,t) + \int\limits_0^t \big(v_{tt} - a^2 \Delta v\big)\, d\tau
  = f(x,t) ,
\]
так как подынтегральное выражение равно нулю при $\tau < t$.
\end{proof}

\paragraph{Явный вид.} Задача для $v$
\[
  \begin{cases}
    v_{tt} - a^2 \Delta v = 0, & x \in \R^3,\ t > \tau, \\[2pt]
    v|_{t=\tau} = 0, \quad v_t|_{t=\tau} = f(x,\tau)
  \end{cases}
\]
— это задача~\eqref{eq:wave3-free} со сдвинутым началом отсчёта времени,
$\varphi \equiv 0$ и $\psi = f(\cdot,\tau)$. Применяя
формулу~\eqref{eq:kirchhoff-free} и интегрируя по $\tau$:
\[
  u(x,t) = \frac{1}{4\pi a^2} \int\limits_0^{t} \frac{1}{t-\tau}
  \int\limits_{S_{x,a(t-\tau)}} f(y,\tau)\, ds\, d\tau .
\]

\subsubsection*{Формула Кирхгофа в общем виде}

Складывая оба вклада, для задачи с данными $\varphi, \psi$ и правой частью $f$:
\[
  \boxed{
  \begin{aligned}
    u(x,t) = \frac{1}{4\pi a^2}\Bigg[ &
      \dd{}{t}\left(\frac1t \int\limits_{S_{x,at}} \varphi(y)\, ds_y\right)
      + \frac1t \int\limits_{S_{x,at}} \psi(y)\, ds_y \\[4pt]
    & + \int\limits_0^{t} \frac{1}{t-\tau}
      \int\limits_{S_{x,a(t-\tau)}} f(y,\tau)\, ds_y\, d\tau
    \Bigg]
  \end{aligned}}
\]

\begin{center}
\textbf{— формула Кирхгофа.}
\end{center}

\inputfig{figures/fig-cone-kirchhoff}

\begin{remark}
В тетради к этому рисунку приписано: \emph{область влияния — это только
образующие конуса, в отличие от одномерного случая}. То есть и терминология
здесь та же, что в \S1 («область влияния \emph{на} точку» = область
зависимости), и подчёркнуто главное отличие от $n=1$: работает не весь конус,
а только его боковая поверхность.
\end{remark}

\begin{idea}[title={Принцип Гюйгенса}]
Во все три интеграла входят только \emph{сферы}, то есть боковая поверхность
конуса, а не его внутренность. Поэтому возмущение, сосредоточенное в момент
$t=0$ в малой окрестности точки $y_0$, проходит через точку $x$ за конечное
время $|x-y_0|/a$ и после этого не оставляет следа — сигнал имеет чёткий задний
фронт.

\medskip
При $n=1$ это не так: в формуле Даламбера интеграл от $\psi$ берётся по всему
отрезку $\Bxat{x}{at}$, и «хвост» остаётся навсегда — это видно по ступеньке в
примере с $\psi = \delta$ из \S1.
\end{idea}

\subsubsection*{Замечание: метод спуска и формула Пуассона}

\begin{remark}
Рассмотрим двумерную задачу
\[
  \begin{cases}
    u_{tt} - a^2 \Delta u = f(x,t), & x \in \R^2,\ t > 0, \\[2pt]
    u|_{t=0} = \varphi(x), \quad u_t|_{t=0} = \psi(x) .
  \end{cases}
\]
Она получается из трёхмерной, если считать все данные не зависящими от третьей
координаты:
\[
  f(x_1,x_2,x_3) = f(x_1,x_2) .
\]
Подстановка в формулу Кирхгофа и интегрирование по $x_3$ даёт \textbf{формулу
Пуассона}.
\end{remark}

\begin{warn}
Сама формула Пуассона не выписана ни на доске, ни в тетради. В тетради есть
только приписка: «получим формулу Пуассона — там будет интеграл по шару в
двумерии, проекции трёхмерной сферы». Вывод, по-видимому, в следующей лекции.
\end{warn}

\begin{idea}
Приём называется \textbf{методом спуска}: задача в $\R^2$ решается как задача в
$\R^3$ с данными, постоянными вдоль $x_3$. Интеграл по сфере $S_{x,at} \subset
\R^3$ при этом проецируется на круг $\Bxat{x}{at} \subset \R^2$, причём сфера
накрывает круг дважды — отсюда в формуле Пуассона появляется интеграл по
\emph{кругу} со знаменателем $\sqrt{a^2t^2 - |y-x|^2}$.

\medskip
Принципиальное следствие: в $\R^2$ интегрирование идёт по кругу, а не по
окружности, значит принцип Гюйгенса \textbf{не выполняется} — как и при $n=1$.
\end{idea}


\newpage
\subsection{Теорема единственности}

\stepcounter{equation}
\begin{equation}\tag{1}\label{eq:uniq-problem}
  \begin{cases}
    \Box u \equiv u_{tt} - a^2 \Delta u = 0, & x \in \R^n,\ t > 0, \\[2pt]
    u|_{t=0} = 0, \quad u_t|_{t=0} = 0, \\[4pt]
    u \in C^2\big(\R^n \times (0,\infty)\big) \cap
         C^{2,1}\big(\R^n \times [0,\infty)\big) .
  \end{cases}
\end{equation}

\begin{warn}
На доске номер параграфа не был проставлен; в тетради он есть — \S4, так что
нумерация верна. Класс гладкости записан как
$C^{2,1}(\R^n \times [0,\infty))$ с пометкой «это условия для классического
решения»: по смыслу две производные по $x$ и одна по $t$, непрерывные вплоть
до границы $t=0$, чтобы начальные условия имели смысл.
\end{warn}

\begin{theorem}[единственность]
Единственное решение задачи~\eqref{eq:uniq-problem} — тождественный нуль:
$u \equiv 0$.
\end{theorem}

\begin{idea}
Это и есть теорема единственности для задачи Коши: если два решения имеют
одинаковые данные $\varphi, \psi$ и одинаковую правую часть $f$, то их разность
решает~\eqref{eq:uniq-problem} и, значит, равна нулю. Заметим, что $n$
произвольно — формул Даламбера и Кирхгофа здесь не требуется.
\end{idea}

\subsubsection*{Энергетическое тождество}

Продифференцируем плотность энергии:
\[
  \frac12 \dd{}{t}\Big(u_t^2 + a^2 |\nabla u|^2\Big)
  = u_t u_{tt} + a^2 \sum_{i=1}^{n} u_{x_i} u_{x_i t}
  = a^2 \sum_{i=1}^{n} \dd{}{x_i}\big(u_{x_i} u_t\big) ,
\]
то есть

\stepcounter{equation}
\begin{equation}\tag{2}\label{eq:energy-identity}
  \frac12 \dd{}{t}\Big(u_t^2 + a^2 |\nabla u|^2\Big)
  - a^2 \sum_{i=1}^{n} \dd{}{x_i}\big(u_{x_i} u_t\big) = 0 .
\end{equation}

\begin{idea}
Последнее равенство в цепочке — это само уравнение. Действительно,
\[
  a^2 \sum_i \dd{}{x_i}\big(u_{x_i} u_t\big)
  = a^2 \underbrace{\Delta u}_{=\,u_{tt}/a^2}\, u_t
    + a^2 \sum_i u_{x_i} u_{x_i t}
  = u_t u_{tt} + a^2 \sum_i u_{x_i} u_{x_i t} .
\]
Тождество~\eqref{eq:energy-identity} — закон сохранения энергии в
дифференциальной форме: $\partial_t(\text{плотность энергии}) +
\operatorname{div}(\text{поток энергии}) = 0$, где поток равен
$-a^2 u_t \nabla u$.
\end{idea}

\subsubsection*{Интегрирование по усечённому конусу}

Зафиксируем точку $(x_0,t_0)$ и возьмём усечённый конус
\[
  D = \big\{\, (x,t) : \ 0 \le t \le t_1, \ \ |x - x_0| \le a(t_0 - t) \,\big\},
  \qquad t_1 \in (0, t_0) .
\]
Его граница состоит из трёх кусков: нижнего основания $\Gamma_0$ (при $t=0$),
верхнего среза $\Gamma_1$ (при $t = t_1$) и боковой поверхности $\Gamma$.

\inputfig{figures/fig-energy-frustum}

Интегрируем~\eqref{eq:energy-identity} по $D$ и применяем формулу
Гаусса—Остроградского в $\R^{n+1}$; через $n = (n_1,\dots,n_n, n_t)$ обозначена
внешняя нормаль:
\begin{align*}
  0 &= \int\limits_D \left[\frac12 \dd{}{t}(\,\cdot\,)
       - a^2 \sum_{i=1}^{n} \dd{}{x_i}(\,\cdot\,)\right] dx\, dt \\[4pt]
    &= \int\limits_{\partial D} \left[\frac12 n_t
       \Big(u_t^2 + a^2 |\nabla u|^2\Big)
       - a^2 \sum_{i=1}^{n} n_i\, u_{x_i} u_t \right] ds
     = \int\limits_{\Gamma_0} \cdot
     + \int\limits_{\Gamma_1} \cdot
     + \int\limits_{\Gamma} \cdot
\end{align*}

Разберём три слагаемых по отдельности.

\paragraph{На $\Gamma_0$.} Нулевые начальные данные дают $u_t|_{t=0} = 0$ и
$\nabla u|_{t=0} = 0$, поэтому
\[
  \int\limits_{\Gamma_0} \cdot \ = 0 .
\]

\paragraph{На $\Gamma_1$.} Здесь нормаль равна $(0,\dots,0,1)$, значит
\[
  \int\limits_{\Gamma_1} \cdot
  = \frac12 \int\limits_{\Gamma_1} \Big(u_t^2 + a^2 |\nabla u|^2\Big) ds
  \ \ge\ 0 .
\]

\paragraph{На $\Gamma$.} Боковая поверхность характеристическая, и для её
нормали выполняется
\[
  n_t^2 = a^2 \sum_{i=1}^{n} n_i^2 .
\]
Умножим и разделим подынтегральное выражение на $n_t$ и воспользуемся этим
соотношением:
\begin{align*}
  \int\limits_{\Gamma} \cdot\ ds
  &= \int\limits_{\Gamma} \frac{1}{n_t}\left[
       \frac12 n_t^2 u_t^2
       + \frac{a^2}{2} \sum_{i=1}^{n} n_t^2 u_{x_i}^2
       - a^2 \sum_{i=1}^{n} n_t n_i\, u_{x_i} u_t
     \right] ds \\[4pt]
  &= \int\limits_{\Gamma} \frac{a^2}{2 n_t} \sum_{i=1}^{n}
     \big(n_i u_t - n_t u_{x_i}\big)^2 ds \ \ge\ 0 .
\end{align*}

\begin{idea}
Фокус ровно в одном месте: $\tfrac12 n_t^2 u_t^2 = \tfrac{a^2}{2}\sum_i n_i^2
u_t^2$ по характеристическому соотношению. После этого все три слагаемых
собираются в полный квадрат
\[
  \frac{a^2}{2 n_t} \sum_i \Big(n_i^2 u_t^2 - 2 n_t n_i u_{x_i} u_t
  + n_t^2 u_{x_i}^2\Big) .
\]
Знак обеспечен тем, что $n_t > 0$: конус раскрыт вниз, и внешняя нормаль к
боковой поверхности имеет положительную компоненту по времени.
\end{idea}

\subsubsection*{Завершение доказательства}

Сумма трёх неотрицательных слагаемых равна нулю, значит каждое из них нулевое. В
частности
\[
  \int\limits_{\Gamma_1} \Big(u_t^2 + a^2 |\nabla u|^2\Big) ds = 0
  \qquad \forall\, t_1 \in [0,t_0] .
\]
Подынтегральное выражение неотрицательно и непрерывно, поэтому внутри конуса
\[
  u_t = 0, \qquad u_{x_i} = 0, \quad i = \overline{1,n}
  \qquad\Longrightarrow\qquad u = \mathrm{const}
  \qquad\Longrightarrow\qquad u \equiv 0 ,
\]
и в частности $u(x_0,t_0) = 0$. Точка $(x_0,t_0)$ была произвольной.
\hfill$\blacksquare$

\begin{idea}
Доказательство попутно даёт то, что в \S1 и \S3 было вычитано из явных формул:
\textbf{конечность скорости распространения}. Всё рассуждение велось внутри
одного конуса, поэтому если данные равны нулю лишь на шаре
$|x - x_0| \le a t_0$, то $u(x_0,t_0) = 0$ всё равно — значение в вершине
определяется только конусом зависимости. Метод при этом не зависит от
размерности и переносится на уравнения с переменными коэффициентами, где
никаких явных формул нет.
\end{idea}


\newpage
\section{Уравнения в частных производных первого порядка}
\subsection{Линейные уравнения}

\begin{warn}
Название главы взято из тетради дословно: «Уравнения в частных производных
первого порядка». У первой главы названия нет ни на доске, ни в тетради —
изложение там начинается сразу с \S1, так что заголовок «Волновое уравнение»
придуман для этого конспекта.
\end{warn}

Однородное уравнение в частных производных первого порядка:

\stepcounter{equation}
\begin{equation}\tag{1}\label{eq:fo-homog}
  \sum_{i=1}^{n} a_i(x)\, u_{x_i} \equiv \bar a(x) \cdot \nabla u = 0 ,
\end{equation}
где коэффициенты $a_i$ гладкие и $\bar a \ne 0$.

\inputfig{figures/fig-char-field}

\subsubsection*{Пример: уравнение переноса}

\begin{example}
Среда (жидкость, газ) движется параллельно оси $x$ с постоянной скоростью $a$.
Пусть $u(x,t)$ — концентрация частиц (числовая плотность, то есть число частиц
на единицу длины).

\medskip
\textbf{Закон сохранения.} Скорость изменения числа частиц на отрезке
$[x_1,x_2]$ равна разности втекающего и вытекающего потоков:
\[
  \dd{}{t} \int\limits_{x_1}^{x_2} u(x,t)\, dx = j(x_1,t) - j(x_2,t),
\]
где $j$ — поток частиц. Для среды, которая просто сносит частицы со скоростью
$a$,
\[
  j(x,t) = a\, u(x,t) .
\]
Подставляя и переписывая правую часть как интеграл от производной:
\[
  \int\limits_{x_1}^{x_2} u_t\, dx + a \int\limits_{x_1}^{x_2} u_x\, dx = 0 .
\]
Отрезок $[x_1,x_2]$ произволен, подынтегральное выражение непрерывно, поэтому
\[
  \boxed{\ u_t + a\, u_x = 0\ }
\]
— линейное однородное \textbf{уравнение переноса}. В обозначениях
задачи~\eqref{eq:fo-homog} направляющий вектор равен $\bar a = (a, 1)$ в
плоскости $(x,t)$.
\end{example}

\inputfig{figures/fig-transport-lines}

В $\R^n$ та же выкладка для среды, движущейся со скоростью $a$ в направлении
единичного вектора $\hat n$, даёт
\[
  u_t + a\, \dd{u}{\hat n} = 0
  \qquad\Longleftrightarrow\qquad
  u_t + \bar a \cdot \nabla u = 0, \qquad \bar a = a\,\hat n ,
\]
с направляющим вектором $(\bar a, 1)$ в пространстве-времени.

\subsubsection*{Геометрическая интерпретация}

Пусть $u$ — гладкое решение уравнения~\eqref{eq:fo-homog}. Тогда:

\begin{itemize}[leftmargin=1.6em]
  \item $u(x) = C$ — поверхность уровня;
  \item $\nabla u$ — нормаль к этой поверхности;
  \item условие $\bar a \cdot \nabla u = 0$ означает $\bar a \perp \nabla u$, то
        есть \textbf{поле направлений $\bar a$ касается поверхностей уровня}.
\end{itemize}

\begin{example}
Уравнение $u_x = 0$ для $u = u(x,y)$: здесь $\bar a = (1,0)$. Поверхности
(линии) уровня — прямые $y = \mathrm{const}$, и общее решение
\[
  u = F(y) ,
\]
то есть произвольная функция от $y$. Линия уровня $F(y) = C$.
\end{example}

\begin{idea}
Отсюда видна общая схема решения уравнения~\eqref{eq:fo-homog}: надо найти
функции, постоянные вдоль интегральных кривых поля $\bar a$ (их называют
\emph{первыми интегралами}), и общее решение будет произвольной функцией от
них. Поэтому следующий шаг — выписать ОДУ для этих кривых.
\end{idea}

\subsubsection*{Характеристическая система}

Интегральные кривые поля $\bar a$ задаются системой ОДУ

\stepcounter{equation}
\begin{equation}\tag{2}\label{eq:char-sys-prop}
  \frac{dx_1}{a_1(x)} = \dots = \frac{dx_n}{a_n(x)} .
\end{equation}
Вводя общее значение этих отношений как $ds$,
\[
  \frac{dx_1}{a_1(x)} = \dots = \frac{dx_n}{a_n(x)} = ds ,
\]
получаем

\stepcounter{equation}
\begin{equation}\tag{3}\label{eq:char-sys}
  \frac{dx_i}{ds} = a_i(x), \qquad i = \overline{1,n} ,
\end{equation}
с начальным условием

\stepcounter{equation}
\begin{equation}\tag{4}\label{eq:char-init}
  x_i(s_0) = x_i^{(0)}, \qquad i = \overline{1,n} ,
\end{equation}
или в векторной форме
\[
  \begin{cases}
    \dfrac{d\vec x}{ds} = \bar a(x), \\[8pt]
    \vec x(s_0) = \vec x_0 .
  \end{cases}
\]

\begin{idea}
Система~\eqref{eq:char-sys} \emph{автономна}: правая часть не зависит от $s$
явно. Это и даёт свойство, сформулированное ниже, — сдвиг параметра вдоль
решения снова даёт решение.
\end{idea}

\subsubsection*{Пример: гармонический осциллятор}

\begin{example}
Уравнение движения груза на пружине:
\[
  \begin{cases}
    m\ddot x = -\kappa x, \\
    x(t_0) = x_0, \quad \dot x(t_0) = v_0 .
  \end{cases}
\]
Сводим к системе первого порядка, вводя скорость $v$:
\[
  \begin{cases}
    \dot x = \dfrac{dx}{dt} = v, \\[6pt]
    m\dot v = m\dfrac{dv}{dt} = -\kappa x, \\[6pt]
    x(t_0) = x_0, \quad v(t_0) = v_0 .
  \end{cases}
\]
В обозначениях задачи~\eqref{eq:char-sys}--\eqref{eq:char-init} это значит
\[
  x_1 = x, \qquad x_2 = v, \qquad s = t ,
\]
а роль поля $\bar a$ играет
\[
  a_1(x,v) = v, \qquad a_2(x,v) = -\frac{\kappa x}{m} = -\omega^2 x,
  \qquad \omega^2 = \frac{\kappa}{m} .
\]
Решение:
\[
  x(t) = \alpha \sin(\omega t + \beta), \qquad
  v(t) = \alpha \omega \cos(\omega t + \beta) ,
\]
а траектории в фазовом пространстве — эллипсы
\[
  \omega^2 x^2 + v^2 = \alpha^2 \omega^2 .
\]
\end{example}

\subsubsection*{Траектории и интегральные кривые}

На этом примере удобно развести два понятия:

\begin{itemize}[leftmargin=1.6em]
  \item $(x(t), v(t)) \in \R^2$ — \textbf{фазовое пространство}, то есть
        пространство зависимых переменных;
  \item $(x(t), v(t), t) \in \R^3$ — \textbf{расширенное фазовое пространство};
  \item $(x(t), v(t), t)$ — \textbf{интегральная кривая} (график решения);
  \item $(x(t), v(t))$ — \textbf{траектория} (в фазовом пространстве).
\end{itemize}

\begin{proposition}
Пусть две интегральные кривые проходят через точки
\[
  (x_0, v_0, t_0) \qquad\text{и}\qquad (x_0, v_0, t_1) .
\]
Тогда они имеют одну и ту же траекторию.
\end{proposition}

В общем случае, для системы~\eqref{eq:char-sys}:

\begin{itemize}[leftmargin=1.6em]
  \item $\vec x(s)$ — траектория;
  \item $(\vec x(s), s) \in \R^{n+1}$ — расширенное фазовое пространство;
  \item $(\vec x(s), s)$ — интегральная кривая (график решения);
  \item интегральные кривые, проходящие через точки $(\vec x_0, s_0)$ и
        $(\vec x_0, s_1)$, имеют одну и ту же траекторию.
\end{itemize}

\begin{idea}
Причина — автономность: если $\vec x(s)$ решает~\eqref{eq:char-sys}, то
$\vec x(s + c)$ тоже решает её при любом $c$. Сдвиг по времени двигает
интегральную кривую вдоль оси $s$, но её проекция на фазовое пространство —
траектория — не меняется. Поэтому траектории не пересекаются и разбивают
фазовое пространство на непересекающиеся кривые, тогда как через одну точку
фазового пространства проходит целое однопараметрическое семейство
интегральных кривых.
\end{idea}

\begin{exercise}
Как выглядят интегральные кривые для гармонического осциллятора?
\end{exercise}

\inputfig{figures/fig-oscillator-phase}

\begin{idea}
Траектория — эллипс в плоскости $(x,v)$, проходимый с постоянной угловой
скоростью $\omega$. Добавляя ось $t$, получаем \textbf{эллиптическую винтовую
линию} в $\R^3$: проекция на $(x,v)$ — тот же эллипс, шаг по $t$ равен периоду
$T = 2\pi/\omega$. Разным начальным моментам $t_0$ отвечают винтовые линии,
сдвинутые друг относительно друга вдоль оси $t$, — ровно то, о чём утверждение
выше. Положению равновесия $x = v = 0$ отвечает вырожденная траектория-точка и
прямая, параллельная оси $t$.
\end{idea}

\subsubsection*{Траектории с общей точкой}

\begin{proposition}\label{prop:traj-coincide}
Траектории, имеющие общую точку, совпадают.
\end{proposition}

\begin{proof}
Пусть $\vec x_1(s)$ — решение~\eqref{eq:char-sys}--\eqref{eq:char-init} с
начальным условием $(\vec x_0, s_1)$, а $\vec x_2(s)$ — решение с начальным
условием $(\vec x_0, s_2)$. Поле $\bar a(x)$ гладкое, поэтому работает теорема
единственности для ОДУ. Утверждается, что
\[
  \vec x_2(s) = \vec x_1(s - s_2 + s_1) .
\]

Правая часть решает~\eqref{eq:char-sys} (система автономна, сдвиг
параметра снова даёт решение), а при $s = s_2$ она равна
$\vec x_1(s_1) = \vec x_0$. Значит она решает ту же задачу Коши, что и
$\vec x_2$, и по единственности обе функции совпадают. Кривые отличаются лишь
сдвигом параметра, то есть совпадают как траектории.
\end{proof}

\begin{idea}
Геометрический смысл: траектория — один и тот же объект, пройденный со сдвигом
параметра, и через каждую точку фазового пространства проходит ровно одна
траектория. Это то же самое утверждение, что и про проекции интегральных
кривых: кривые, сдвинутые вдоль оси $s$, проектируются в одну траекторию.
\end{idea}

\begin{definition}
Траектория системы~\eqref{eq:char-sys} называется её
\textbf{характеристикой}. Характеристиками уравнения~\eqref{eq:fo-homog}
называются характеристики системы~\eqref{eq:char-sys}.
\end{definition}

\begin{remark}
Характеристики системы~\eqref{eq:char-sys} описываются характеристической
системой~\eqref{eq:char-sys-prop}, в которой параметра $s$ уже нет: он нужен
был только чтобы перейти от пропорции к автономной системе, а сама кривая от
параметризации не зависит.
\end{remark}

\begin{warn}
Смазанная фраза на доске («хар-ки системы (3) описываются характеристическими
системами») — это именно замечание выше: в тетради она записана как «хар-ки
системы (3) описываются хар. системой (2) (нет переменной $s$)».
\end{warn}

\subsubsection*{Геометрическая интерпретация характеристик}

\begin{itemize}[leftmargin=1.6em]
  \item поле $\bar a$ в каждой точке касается характеристики;
  \item обратно, кривая, в каждой точке которой поле $\bar a$ касательно,
        является характеристикой.
\end{itemize}

Оба утверждения выводятся из пропорции~\eqref{eq:char-sys-prop} и теоремы
единственности для ОДУ: касание поля в каждой точке — это и есть
покоординатная пропорциональность $dx_i$ и $a_i(x)$, а дальше через каждую
точку проходит ровно одна такая кривая.

Пусть $u(\vec x)$ — гладкая функция. Тогда вдоль характеристики

\stepcounter{equation}
\begin{equation}\tag{5}\label{eq:char-deriv}
  \frac{d\,u(\vec x(s))}{ds}
  = \sum_{i=1}^{n} u_{x_i} \frac{dx_i}{ds}
  = \sum_{i=1}^{n} a_i(x)\, u_{x_i}
  = \bar a(x) \cdot \nabla u .
\end{equation}

\begin{remark}
Интерпретация соотношения~\eqref{eq:char-deriv}: \textbf{на характеристиках
оператор в частных производных $\bar a \cdot \nabla$ превращается в обыкновенную
производную по $s$}. Это и есть весь метод характеристик.
\end{remark}

\subsubsection*{Первые интегралы}

\begin{definition}
Функция $u(\vec x) \not\equiv \mathrm{const}$ называется \textbf{первым
интегралом} системы~\eqref{eq:char-sys}, если она постоянна на любом решении
этой системы.
\end{definition}

\begin{proposition}
Функция $u(x)$ является решением уравнения~\eqref{eq:fo-homog} тогда и только
тогда, когда она является первым интегралом системы~\eqref{eq:char-sys}.
\end{proposition}

\begin{proof}
Немедленно из~\eqref{eq:char-deriv}: равенство $\bar a \cdot \nabla u = 0$ в
точности означает, что $\frac{d}{ds} u(\vec x(s)) = 0$ вдоль каждой
характеристики.
\end{proof}

\begin{definition}
Первые интегралы $u_1,\dots,u_k$ называются \textbf{независимыми}, если
\[
  \operatorname{rank}
  \frac{\partial(u_1,\dots,u_k)}{\partial(x_1,\dots,x_n)} = k ,
\]
то есть не существует $F \not\equiv 0$, для которой
$F(u_1,\dots,u_k) \equiv 0$ в окрестности точки.
\end{definition}

\begin{warn}
Знаменатель якобиана: на доске он смазан, в тетради записан правильно —
$\partial(x_1,\dots,x_n)$, с пометкой «ранг максимален».

\medskip
А вот словесная расшифровка независимости и на доске, и в тетради записана
криво: «не существует функции $F(u_1,\dots,u_k) \not\equiv \mathrm{const}$ в
окрестности точки». По смыслу имеется в виду «не существует нетривиальной
гладкой $F$, для которой $F(u_1,\dots,u_k) \equiv \mathrm{const}$», то есть
функциональная независимость. В тетради рядом стоит более внятная приписка:
\emph{все градиенты не коллинеарны} — это ровно условие на ранг.
\end{warn}

\begin{theorem}\label{thm:first-integrals}
Система~\eqref{eq:char-sys} имеет ровно $n-1$ независимых первых интегралов;
обозначим их через $u_1,\dots,u_{n-1}$. Всякий первый интеграл имеет вид
\[
  F(u_1,\dots,u_{n-1}), \qquad F \not\equiv \mathrm{const} ,
\]
и обратно: для любой гладкой $F \not\equiv \mathrm{const}$ функция
$F(u_1,\dots,u_{n-1})$ является первым интегралом.
\end{theorem}

\begin{warn}
На доске обе половины утверждения остались недописанными; тетрадь их
подтверждает дословно: «всякий первый интеграл имеет вид
$F(u_1,\dots,u_{n-1})$, где $F$ гладкая и $F \not\equiv \mathrm{const}$, и
наоборот — всякая такая гладкая функция есть первый интеграл». Так что
восстановление верное.
\end{warn}

\begin{corollary*}
Общее решение уравнения~\eqref{eq:fo-homog} имеет вид
$u = F(u_1,\dots,u_{n-1})$.
\end{corollary*}

\begin{idea}
Почему именно $n-1$: траектории заполняют $\R^n$ однопараметрическими кривыми,
так что «поперёк» них остаётся $n-1$ независимых направлений. Первый интеграл —
это функция, постоянная вдоль траекторий, то есть функция на множестве самих
траекторий, а это множество $(n-1)$-мерно.
\end{idea}

\subsubsection*{Как искать первые интегралы}

Решаем систему~\eqref{eq:char-sys-prop} и получаем $n-1$ соотношений — это
решение в неявном виде:
\[
  \begin{cases}
    W_1(x,\, C_1,\dots,C_{n-1}) = 0, \\
    \ \ \vdots \\
    W_{n-1}(x,\, C_1,\dots,C_{n-1}) = 0,
  \end{cases}
\]
после чего разрешаем их относительно констант:
\[
  \begin{cases}
    u_1(x) = C_1, \\
    \ \ \vdots \\
    u_{n-1}(x) = C_{n-1} .
  \end{cases}
\]
Слева в этих равенствах и стоят искомые первые интегралы, а сама система при
произвольных $C_i$ задаёт \emph{произвольную} траекторию
системы~\eqref{eq:char-sys}.

\medskip
Чтобы получить \emph{конкретную} траекторию — ту, что проходит через точку
$x_0$, то есть решение задачи~\eqref{eq:char-sys}--\eqref{eq:char-init}, —
надо подставить вместо констант значения в этой точке:
\[
  \begin{cases}
    u_1(x) = u_1(x_0), \\
    \ \ \vdots \\
    u_{n-1}(x) = u_{n-1}(x_0) .
  \end{cases}
\]

\begin{idea}[title={Почему можно просто подставить $x_0$}]
Пометка «почему?» в тетради отсылает к системе~\eqref{eq:char-sys} с начальным
условием~\eqref{eq:char-init}. Вдоль траектории функции $u_i$ по определению
не меняются, а точка $x_0$ на этой траектории лежит. Значит константы обязаны
равняться $C_i = u_i(x_0)$ — другого выбора нет.
\end{idea}

\subsubsection*{Пример: уравнение переноса}

\begin{example}
Для $u = u(x,t)$ и уравнения $u_t + a u_x = 0$ характеристическая
система~\eqref{eq:char-sys-prop} выглядит как
\[
  \frac{dx}{a} = \frac{dt}{1} ,
\]
откуда
\[
  x - at = C \quad\text{— характеристики}, \qquad
  x - at \quad\text{— первый интеграл} .
\]
Общее решение:
\[
  u(x,t) = F(x - at) ,
\]
направляющий вектор $\bar a = (a,1)$ (см. рис.~\ref{fig:transport-lines}).

\medskip
Добавим начальное условие $u|_{t=0} = \varphi(x)$. Тогда $F(x) = \varphi(x)$, и
решение задачи равно
\[
  \boxed{\ u(x,t) = \varphi(x - at)\ }
\]
\end{example}

\begin{warn}
В тетради последняя строка не дописана: «$\Rightarrow$ реш-е задачи
$u(x,t) =$ \textbf{???}». Ответ восстановлен подстановкой $F = \varphi$.
\end{warn}

\begin{exercise}
Найти характеристики и общее решение уравнения
\[
  u_t + a u_x + b u_y = 0, \qquad u = u(x,y,t),
  \qquad \bar a = (a,\, b,\, 1) .
\]
\end{exercise}

\begin{idea}
Характеристическая система: $\dfrac{dx}{a} = \dfrac{dy}{b} = \dfrac{dt}{1}$.
Независимых первых интегралов здесь $n-1 = 2$ штуки: $x - at$ и $y - bt$.
Характеристики — прямые $x - at = C_1$, $y - bt = C_2$, а общее решение
$u = F(x - at,\ y - bt)$.
\end{idea}

\subsubsection*{Задача Коши на кривой}

Пусть данные заданы не на прямой $t=0$, а на кривой
\[
  \gamma = \big\{\, (\xi(\tau),\, \eta(\tau)) \,\big\},
  \qquad u|_\gamma = \varphi(\tau) .
\]

\begin{definition}
\textbf{Условие трансверсальности}: поле направлений $\bar a$ нигде не касается
$\gamma$.
\end{definition}

\inputfig{figures/fig-cauchy-gamma}

Подставляя общее решение $u = F(x - at)$ в начальное условие, получаем уравнение
на $F$:

\stepcounter{equation}
\begin{equation}\tag{6}\label{eq:transversal-eq}
  F\big(\xi(\tau) - a\,\eta(\tau)\big) = \varphi(\tau) .
\end{equation}

\begin{exercise}[теорема об обратной функции]
Докажите, что если $\xi, \eta, \varphi$ гладкие и выполнено условие
трансверсальности, то уравнение~\eqref{eq:transversal-eq} можно однозначно
разрешить относительно $F$, и значит решение задачи Коши существует и
единственно в окрестности $\gamma$.

\medskip
\emph{Подсказка:} рассмотреть $g(\tau) = \xi(\tau) - a\,\eta(\tau)$.
\end{exercise}

\begin{idea}
Разбор подсказки. Касательный вектор к $\gamma$ равен $(\xi', \eta')$, а
направляющий вектор характеристики — $\bar a = (a, 1)$. Они параллельны тогда и
только тогда, когда
\[
  \det \begin{pmatrix} \xi' & a \\ \eta' & 1 \end{pmatrix}
  = \xi'(\tau) - a\,\eta'(\tau) = 0 .
\]
Но это в точности $g'(\tau)$. Значит условие трансверсальности равносильно
$g'(\tau) \ne 0$: функция $g$ строго монотонна и потому локально обратима.
Тогда~\eqref{eq:transversal-eq} даёт
\[
  F = \varphi \circ g^{-1} ,
\]
то есть $F$ восстанавливается однозначно, а решение равно
$u(x,t) = \varphi\big(g^{-1}(x - at)\big)$. Гладкость $F$ наследуется от
$\varphi$ и $g^{-1}$.

\medskip
Короче говоря: \textbf{трансверсальность $\iff$ якобиан отличен от нуля}, и всё
упражнение — это теорема об обратной функции.
\end{idea}

\subsubsection*{Что бывает без трансверсальности}

\begin{example}
\[
  \begin{cases}
    u_x = 0, \\
    u|_{y = x^3} = x^2 .
  \end{cases}
\]
Здесь $\bar a = (1,0)$, характеристики — прямые $y = \mathrm{const}$, а кривая
$y = x^3$ касается одной из них в начале координат: её касательная
$(1, 3x^2)$ при $x = 0$ равна $(1,0)$.

\medskip
Общее решение $u = F(y)$; из начального условия $F(x^3) = x^2$, то есть
\[
  u = y^{2/3} ,
\]
и это функция, \textbf{негладкая на оси} $Ox$.
\end{example}

\inputfig{figures/fig-tangency}

\begin{idea}
Трансверсальность нарушена ровно в одной точке — и этого уже достаточно, чтобы
решение потеряло гладкость. Причина видна по формуле $F = \varphi \circ g^{-1}$
из предыдущего пункта: в точке касания $g' = 0$, обратная функция $g^{-1}$
негладкая, и эта негладкость переносится на $F$. Здесь $g(x) = x^3$,
$g^{-1}(y) = y^{1/3}$, $\varphi(x) = x^2$ — и получается ровно $y^{2/3}$.
\end{idea}


\newpage
\subsection{Полулинейные уравнения}

\stepcounter{equation}
\begin{equation}\tag{1}\label{eq:semilin}
  \sum_{i} a_i(x)\, u_{x_i} \equiv \bar a \cdot \nabla u = f(x, u) .
\end{equation}

\begin{warn}
Номер параграфа не проставлен ни на доске, ни в тетради: заголовок идёт сразу
после контрпримера с касанием. Здесь он взят следующим по порядку. В тетради
он записан как «Полулинейное ур-е ЧП 1-го \emph{рода}» — по смыслу речь, как и
везде, о первом \emph{порядке}. Нумерация формул в этом параграфе начинается
заново.
\end{warn}

\begin{idea}
Отличие от~\eqref{eq:fo-homog} только в правой части, причём $f$ зависит от $u$,
но \emph{не} от производных — поэтому «полулинейное». Левая часть по-прежнему
линейна по $\nabla u$, значит характеристики те же самые: они определяются одним
лишь полем $\bar a(x)$ и от $f$ не зависят. Меняется лишь то, что вдоль
характеристики $u$ теперь не постоянна, а подчиняется ОДУ.
\end{idea}

\subsubsection*{Задача Коши на поверхности}

Пусть задана поверхность размерности $n-1$:
\[
  S = \Big\{\, \vec\xi(\vec\tau) \equiv
      \big(\xi_1(\vec\tau),\dots,\xi_n(\vec\tau)\big), \quad
      \vec\tau = (\tau_1,\dots,\tau_{n-1}) \,\Big\} ,
\]
и начальное условие
\[
  u|_S = \varphi(\vec\tau) .
\]

\begin{definition}
\textbf{Условие трансверсальности}: поле направлений (то есть характеристики)
нигде не касается $S$.
\end{definition}

\inputfig{figures/fig-surface-transversal}

\subsubsection*{Замена переменных вдоль характеристик}

Выпустим из каждой точки поверхности характеристику:
\[
  \begin{cases}
    \dfrac{d\vec x}{ds} = \bar a(\vec x), \\[8pt]
    \vec x(0) = \vec\xi(\vec\tau) .
  \end{cases}
\]
Это задаёт замену переменных $(s, \vec\tau) \mapsto \vec x$, при которой сама
поверхность получает уравнение
\[
  S = \{\, s = 0 \,\} .
\]

Якобиан этой замены равен
\[
  \begin{pmatrix}
    x_{1s} & x_{1\tau_1} & \dots & x_{1\tau_{n-1}} \\
    \vdots & \vdots      &       & \vdots \\
    x_{ns} & x_{n\tau_1} & \dots & x_{n\tau_{n-1}}
  \end{pmatrix}
  \Bigg|_{s=0}
  = \Big(\, \bar a\big(\vec\xi(\vec\tau)\big) \ \Big|\
    \vec\xi_{\tau_1}(\vec\tau) \ \Big|\ \dots \ \Big|\
    \vec\xi_{\tau_{n-1}}(\vec\tau) \,\Big) \ \ne\ 0 ,
\]
и отличие определителя от нуля — это в точности условие трансверсальности.

\begin{idea}[title={Что означает этот якобиан}]
При $s = 0$ первый столбец равен
\[
  \dd{\vec x}{s}\bigg|_{s=0} = \bar a\big(\vec\xi(\vec\tau)\big) ,
\]
то есть направляющему вектору характеристики, а остальные $n-1$ столбцов —
это $\partial \vec\xi / \partial \tau_j$, то есть базис касательной плоскости к
$S$. Определитель отличен от нуля тогда и только тогда, когда $\bar a$ не лежит
в касательной плоскости, — а это и есть условие трансверсальности.

\medskip
Значит по теореме об обратной функции замена $(s,\vec\tau) \leftrightarrow
\vec x$ локально взаимно однозначна и гладко обратима в окрестности $S$. Именно
здесь трансверсальность и нужна; без неё характеристики «ложатся» на $S$,
замена вырождается, и получается ситуация вроде контрпримера с $y = x^3$
из \S1.
\end{idea}

\begin{warn}
В тетради справа от матрицы стояло просто «$= \bar a(\vec\xi)$» — отсюда и
возникал вопрос. Полная запись есть на другой странице тетради и приведена
выше: первый столбец равен $\bar a(\vec\xi)$, остальные — производные
$\vec\xi$ по $\tau_j$, а всё вместе отлично от нуля по условию
трансверсальности.
\end{warn}

\subsubsection*{Сведение к ОДУ}

Итак, обратное отображение определено, и переменные переходят друг в друга:
\[
  \vec x \ \longleftrightarrow\ (s, \vec\tau) .
\]
Сделаем замену
\[
  u(\vec x) = v(s, \vec\tau) .
\]
По соотношению~\eqref{eq:char-deriv} производная по $s$ — это в точности
$\bar a \cdot \nabla u$, поэтому задача~\eqref{eq:semilin} превращается в
\[
  \begin{cases}
    \dfrac{\partial v}{\partial s}
      \overset{\eqref{eq:char-deriv}}{=} \bar a \cdot \nabla u
      = f(x, u) = \widetilde f(s, \vec\tau, v), \\[10pt]
    v|_{s=0} = \varphi(\vec\tau) .
  \end{cases}
\]
А это обыкновенное дифференциальное уравнение — решать умеем.

\begin{idea}
Это обыкновенное дифференциальное уравнение по $s$, в котором $\vec\tau$ —
просто набор параметров. То есть метод характеристик целиком: уравнение в
частных производных распадается на семейство задач Коши для ОДУ, по одной на
каждую характеристику, а начальные данные для них считываются с поверхности
$S$.

\medskip
Линейный однородный случай \S1 — частный: там $f \equiv 0$, уравнение
принимает вид $\partial v / \partial s = 0$, то есть $v$ постоянна вдоль
характеристики, и решение сводится к переносу данных с $S$.
\end{idea}


\end{document}
